GPU work is organised into thread blocks; threads within a block can
cooperate through shared memory. A warp is a group of 32 threads scheduled
together. Blocks run on streaming multiprocessors (SMs). Registers, shared memory
and block limits constrain concurrent work \cite{cuda_programming_guide}.
Occupancy is the fraction of an SM's supported warps that are active.
Nearby addresses accessed by a warp can share memory transactions
(coalesced access); large per-block tables can reduce concurrent work.
Within this hierarchy, threads can process separate queries, score
different candidates for the same query, or share the coordinate
calculations for one candidate.

Candidate selection keeps the smallest scores from the scanned pool.
Faiss uses warp-level selection \cite{johnson_gpu}; other GPU top-$k$
methods distribute selection over larger arrays \cite{gpu_topk}.
Chapter 5 explains the selection and scheduling used by JHQ-GPU.

